% ============================================================
%  PAUTA PEP 1 — FISICA I PARA INGENIERIA (FORMA A) — USACH
%  Usach Premium
%  Compila con: pdflatex pauta_pep1_forma_a.tex  (2 veces)
% ============================================================

\documentclass[letterpaper, 11pt]{article}

\usepackage[utf8]{inputenc}
\usepackage[spanish]{babel}
\usepackage{amsmath, amssymb}
\usepackage{geometry}
\usepackage{fancyhdr}
\usepackage{enumitem}
\usepackage{graphicx}
\usepackage{hyperref}
\usepackage{parskip}
\usepackage{xcolor}
\usepackage{tcolorbox}
\usepackage{eso-pic}
\usepackage{tikz}
\usetikzlibrary{patterns, arrows.meta, decorations.pathmorphing, calc, babel}
\tcbuselibrary{skins, breakable}

\geometry{letterpaper, margin=2cm, top=3cm, bottom=2.5cm, headheight=2cm}

\definecolor{celesteSuave}{HTML}{F0F7FF}
\definecolor{azulFuerte}{HTML}{1A5276}
\definecolor{verdePaso}{HTML}{1E8449}
\definecolor{rojoAlerta}{HTML}{C0392B}
\definecolor{naranjaAcento}{HTML}{D35400}
\definecolor{enlace}{HTML}{1A5276}

\newcommand{\logoup}{./images/logo_up.png}
\newcommand{\logousach}{./images/logo_usach.png}
\newcommand{\logofondo}{./images/logo_up.png}

\newcommand{\institucion}{Usach Premium}
\newcommand{\curso}{Fisica I para Ingenieria}
\newcommand{\guiasnumero}{Pauta PEP 1 --- Forma A}
\newcommand{\semestre}{1er. Semestre de 2026}
\newcommand{\preparadopor}{Usach Premium.}
\newcommand{\webinstitucion}{www.usachpremium.cl}
\newcommand{\instagraminstitucion}{https://www.instagram.com/usach.premium}
\newcommand{\transparencia}{0.04}
\newcommand{\fraseportada}{"Por y para estudiantes"}

\hypersetup{colorlinks=true, linkcolor=enlace, urlcolor=enlace,
            citecolor=enlace, pdfborder={0 0 0}}

\pagestyle{fancy}
\fancyhf{}
\fancyhead[L]{\includegraphics[height=1.2cm]{\logoup}}
\fancyhead[R]{%
    \begin{minipage}[b]{0.45\textwidth}
        \raggedleft\fontsize{9pt}{11pt}\selectfont
        \textbf{\color{azulFuerte}\institucion} \\
        \curso\ --- \guiasnumero \\
        \textit{\color{gray}@usach.premium}
    \end{minipage}\hspace{0.1cm}%
    \includegraphics[height=1.2cm]{\logousach}}
\fancyfoot[L]{\fontsize{9pt}{11pt}\selectfont \textit{\fraseportada}}
\fancyfoot[R]{\fontsize{9pt}{11pt}\selectfont P\'agina \thepage}
\renewcommand{\headrulewidth}{0.4pt}
\renewcommand{\footrulewidth}{0.4pt}
\fancypagestyle{plain}{\fancyhf{}\renewcommand{\headrulewidth}{0pt}
                                   \renewcommand{\footrulewidth}{0pt}}

\newcommand{\MarcaAgua}{
    \AddToShipoutPicture{
        \AtPageCenter{
            \begin{tikzpicture}[remember picture, overlay]
                \node[opacity=\transparencia] at (0,0)
                    {\includegraphics[width=0.7\textwidth]{\logofondo}};
            \end{tikzpicture}}}}

\newtcolorbox{paso}[2][]{
    colback=celesteSuave, colframe=azulFuerte, fonttitle=\bfseries,
    coltitle=white, colbacktitle=azulFuerte,
    title={\large Paso #2},
    arc=6pt, boxrule=1pt, enhanced, breakable, #1
}
\newtcolorbox{justif}[1]{
    colback=naranjaAcento!8, colframe=naranjaAcento, fonttitle=\bfseries\small,
    title={#1}, arc=5pt, boxrule=0.8pt, enhanced, breakable
}
\newtcolorbox{enunciadobox}[1]{
    colback=azulFuerte!10, colframe=azulFuerte, fonttitle=\bfseries,
    coltitle=white, colbacktitle=azulFuerte, title={#1},
    arc=7pt, boxrule=1.2pt, enhanced, breakable
}
\newtcolorbox{resultado}{
    colback=azulFuerte, colframe=azulFuerte, colupper=white,
    arc=10pt, boxrule=1.5pt, enhanced, drop shadow={black!15}
}

\begin{document}
\thispagestyle{plain}
\MarcaAgua

\AddToShipoutPicture*{%
    \put(54,700){\includegraphics[width=2cm]{\logoup}}
    \put(420,706){\includegraphics[width=5cm]{\logousach}}}

\vspace*{0.5cm}
\begin{center}
    \color{azulFuerte}
    {\huge \textbf{PAUTA DE CORRECCI\'ON}} \\[0.5cm]
    {\Huge \textbf{PEP 1 --- \curso}} \\[0.4cm]
    {\Large Forma A \quad$\bullet$\quad C\'od: 10103-10140} \\[1.2cm]
    \begin{tcolorbox}[colback=celesteSuave, colframe=azulFuerte, arc=10pt,
        width=0.88\textwidth, center, boxrule=1.5pt]
        \centering\vspace{0.2cm}
        {\large \textbf{Contenidos evaluados}}\\[0.4cm]
        \color{black}
        \begin{itemize}[leftmargin=2.5cm, itemsep=3pt]
            \item[\color{azulFuerte}$\Rightarrow$] Cinem\'atica: gr\'afico $v$-$t$ y aceleraci\'on
            \item[\color{azulFuerte}$\Rightarrow$] Ecuaciones de posici\'on y encuentro de m\'oviles
            \item[\color{azulFuerte}$\Rightarrow$] Lanzamiento vertical por etapas (MRUA)
            \item[\color{azulFuerte}$\Rightarrow$] Suma de vectores: componentes, magnitud y direcci\'on
        \end{itemize}
        \vspace{0.15cm}
    \end{tcolorbox}
\end{center}

\vspace{1.0cm}
\begin{center}
    {\Large \textbf{\semestre}} \\[0.6cm]
    {\large \textbf{\preparadopor}}\\[0.4cm]
    {\footnotesize\color{gray} Considera $g = 9{,}8\ \mathrm{m/s^2}$. Resultados finales aproximados a 2 decimales.}
\end{center}

\vspace{1.2cm}
\begin{center}
    {\color{azulFuerte}\hrule height 0.5pt}\vspace{0.3cm}
    \begin{minipage}{0.45\textwidth}\centering
        \textbf{Instagram:} \\\small\href{\instagraminstitucion}{@usach.premium}
    \end{minipage}
    \begin{minipage}{0.45\textwidth}\centering
        \textbf{Web:} \\\small\href{http://\webinstitucion}{\webinstitucion}
    \end{minipage}
\end{center}

\pagenumbering{arabic}
\newpage

% ============================================================
% ============================  P1  ==========================
% ============================================================
\begin{enunciadobox}{Problema 1 --- Cami\'on de reparto aut\'onomo \hfill (20 puntos)}
Cami\'on de masa $M=2195\,\mathrm{kg}$ que se mueve seg\'un el gr\'afico $v$-$t$
(tres etapas). Un vehículo de apoyo ($m=1250\,\mathrm{kg}$) parte del mismo
origen con retraso de $3\,\mathrm{s}$, rapidez inicial $v_0$ desconocida y
aceleraci\'on $\vec{a}=(1{,}5\,\mathrm{m/s^2})\,\hat{\imath}$. Ambos se encuentran
en la \textbf{segunda etapa} y en ese instante el apoyo alcanza $v_f=28\,\mathrm{m/s}$.
\end{enunciadobox}

\medskip
\begin{center}
\begin{tikzpicture}[x=0.28cm,y=0.13cm]
  \draw[-{Latex[length=2mm]},line width=0.6pt] (0,0) -- (30,0) node[below left=1pt]{$t\,(\mathrm{s})$};
  \draw[-{Latex[length=2mm]},line width=0.6pt] (0,0) -- (0,32) node[right=2pt]{$v\,(\mathrm{m/s})$};
  \draw[dashed] (0,24) -- (26,24);
  \draw[dashed] (8,0) -- (8,24);
  \draw[dashed] (20,0) -- (20,24);
  \draw[line width=1pt, azulFuerte] (0,0) -- (8,24) -- (20,24) -- (26,0);
  \node[below] at (8,0) {$8$}; \node[below] at (20,0) {$20$}; \node[below] at (26,0) {$26$};
  \node[left] at (0,24) {$24$}; \node[below left] at (0,0){$0$};
  \node[verdePaso] at (4,13) {\scriptsize Etapa 1};
  \node[verdePaso] at (14,26) {\scriptsize Etapa 2};
  \node[verdePaso] at (24,13) {\scriptsize Etapa 3};
\end{tikzpicture}
\end{center}

% ------------------- P1 (a) -------------------
\begin{paso}{1.a: Aceleraciones $a_1$ y $a_3$ (pendiente del gr\'afico)}
La aceleraci\'on es la pendiente de la recta $v$-$t$ en cada tramo.

\textbf{Etapa 1} ($0\le t\le 8$): la velocidad pasa de $0$ a $24\,\mathrm{m/s}$:
\[
a_1=\frac{\Delta v}{\Delta t}=\frac{24-0}{8-0}=3\ \mathrm{m/s^2}.
\]
\textbf{Etapa 3} ($20\le t\le 26$): la velocidad pasa de $24$ a $0\,\mathrm{m/s}$:
\[
a_3=\frac{0-24}{26-20}=\frac{-24}{6}=-4\ \mathrm{m/s^2}.
\]
\begin{resultado}\centering
$\displaystyle a_1 = 3{,}00\ \mathrm{m/s^2}\qquad\text{y}\qquad a_3 = -4{,}00\ \mathrm{m/s^2}$
\end{resultado}
\end{paso}

% ------------------- P1 (b) -------------------
\begin{paso}{1.b: Ecuaciones de posici\'on $x(t)$ (origen en la posici\'on inicial del cami\'on)}
Se usa $x(t)=x_0+v_0 t+\tfrac12 a t^2$ en cada tramo, encadenando posici\'on y
velocidad al final de cada etapa (la etapa 2 es MRU con $v=24\,\mathrm{m/s}$).

\textbf{Etapa 1} ($0\le t\le 8$): parte del reposo en el origen, $a_1=3$:
\[
x_1(t)=\tfrac12(3)t^2=1{,}5\,t^2,\qquad x_1(8)=96\ \mathrm{m}.
\]
\textbf{Etapa 2} ($8\le t\le 20$): parte de $96\,\mathrm{m}$ con $v=24$ constante:
\[
x_2(t)=96+24\,(t-8),\qquad x_2(20)=384\ \mathrm{m}.
\]
\textbf{Etapa 3} ($20\le t\le 26$): parte de $384\,\mathrm{m}$ con $v=24$ y $a_3=-4$:
\[
x_3(t)=384+24\,(t-20)-2\,(t-20)^2,\qquad x_3(26)=456\ \mathrm{m}.
\]
\begin{resultado}\centering
$\displaystyle x_1(t)=1{,}5\,t^2\quad;\quad x_2(t)=96+24(t-8)\quad;\quad x_3(t)=384+24(t-20)-2(t-20)^2$
\end{resultado}
\end{paso}

% ------------------- P1 (c) -------------------
\begin{paso}{1.c: Rapidez inicial $v_0$ del apoyo y posici\'on de encuentro $x^{*}$}
El apoyo parte en $t=3\,\mathrm{s}$. Con $\tau=t-3$ (su tiempo de movimiento) y $a=1{,}5$:
\[
v_{\text{apoyo}}(\tau)=v_0+1{,}5\,\tau,\qquad
x_{\text{apoyo}}(\tau)=v_0\,\tau+\tfrac12(1{,}5)\tau^2=v_0\,\tau+0{,}75\,\tau^2 .
\]

\textbf{Condici\'on 1 (rapidez en el encuentro).} El apoyo alcanza $v_f=28\,\mathrm{m/s}$:
\[
v_0+1{,}5\,\tau=28 \;\Rightarrow\; v_0=28-1{,}5\,\tau. \tag{i}
\]

\textbf{Condici\'on 2 (misma posici\'on en la etapa 2).} El encuentro ocurre en la
etapa 2 del cami\'on, donde $x_{\text{cam}}=96+24\,(t-8)=96+24(\tau-5)$:
\[
v_0\tau+0{,}75\,\tau^2 = 96+24(\tau-5). \tag{ii}
\]

Reemplazando (i) en (ii):
\[
(28-1{,}5\tau)\tau+0{,}75\tau^2 = 24\tau-24
\;\Longrightarrow\; 0{,}75\,\tau^2-4\tau-24=0
\;\Longleftrightarrow\; 3\tau^2-16\tau-96=0.
\]
\[
\tau=\frac{16+\sqrt{256+1152}}{6}=\frac{16+\sqrt{1408}}{6}\approx 8{,}92\ \mathrm{s}
\;\Rightarrow\; t^{*}=\tau+3\approx 11{,}92\ \mathrm{s}.
\]
\begin{justif}{Verificaci\'on de la etapa}
$t^{*}\approx 11{,}92\,\mathrm{s}\in[8,20]$: el encuentro efectivamente ocurre en la
\textbf{segunda etapa}, como pide el enunciado.
\end{justif}

Rapidez inicial y posici\'on de encuentro:
\[
v_0=28-1{,}5(8{,}92)\approx 14{,}62\ \mathrm{m/s},\qquad
x^{*}=96+24(11{,}92-8)\approx 190{,}09\ \mathrm{m}.
\]
\begin{resultado}\centering
$\displaystyle v_0 \approx 14{,}62\ \mathrm{m/s}\qquad\text{y}\qquad x^{*}\approx 190{,}09\ \mathrm{m}$
\end{resultado}
\end{paso}

\newpage
% ============================================================
% ============================  P2  ==========================
% ============================================================
\begin{enunciadobox}{Problema 2 --- Cohete de prueba (4 etapas) \hfill (20 puntos)}
Lanzamiento vertical desde el suelo. \textbf{Etapa I:} sube libre con
$V_{I0}=80\,\mathrm{m/s}$ (solo gravedad) hasta detenerse. \textbf{Etapa II:}
motores con $a_{II}=4\,\mathrm{m/s^2}$ hasta $h=1000\,\mathrm{m}$. \textbf{Etapa III:}
solo gravedad hasta la altura m\'axima. \textbf{Etapa IV:} frenado uniforme desde
$H=200\,\mathrm{m}$ hasta el reposo al tocar el suelo. Se toma $\hat{\jmath}$ hacia arriba.
\end{enunciadobox}

% ------------------- P2 (a) -------------------
\begin{paso}{2.a: Altura m\'axima $y_{\max}$}
\textbf{Etapa I} (sube y se detiene, $a=-g$):
\[
0=V_{I0}^2-2g\,y_I \;\Rightarrow\; y_I=\frac{80^2}{2(9{,}8)}=326{,}53\ \mathrm{m}.
\]
\textbf{Etapa II} (de $y_I=326{,}53$ con $v=0$ hasta $h=1000$, $a_{II}=4$):
\[
v_{II}^2=0+2(4)(1000-326{,}53)=5387{,}76 \;\Rightarrow\; v_{II}\approx 73{,}40\ \mathrm{m/s}.
\]
\textbf{Etapa III} (solo gravedad, sube hasta detenerse):
\[
y_{\max}=h+\frac{v_{II}^2}{2g}=1000+\frac{5387{,}76}{2(9{,}8)}=1000+274{,}89 .
\]
\begin{resultado}\centering
$\displaystyle y_{\max}\approx 1274{,}89\ \mathrm{m}$
\end{resultado}
\end{paso}

% ------------------- P2 (b) -------------------
\begin{paso}{2.b: Velocidad $\vec{v}_{IV}$ al inicio de la etapa final}
Usando el dato $y_{\max}=1274{,}89\,\mathrm{m}$, el cohete \textbf{cae libremente}
desde el reposo en $y_{\max}$ hasta $H=200\,\mathrm{m}$ (solo gravedad):
\[
v_{IV}^2=2g\,(y_{\max}-H)=2(9{,}8)(1274{,}89-200)=21067{,}84 .
\]
\[
v_{IV}=\sqrt{21067{,}84}\approx 145{,}15\ \mathrm{m/s}\ \ (\text{hacia abajo}).
\]
\begin{justif}{Forma anal\'itica (vector)}
El movimiento es descendente, por lo que la velocidad apunta en $-\hat{\jmath}$.
\end{justif}
\begin{resultado}\centering
$\displaystyle \vec{v}_{IV}\approx -145{,}15\,\hat{\jmath}\ \ \mathrm{m/s}$
\end{resultado}
\end{paso}

% ------------------- P2 (c) -------------------
\begin{paso}{2.c: Aceleraci\'on $\vec{a}_{\mathrm{final}}$ en el aterrizaje}
Usando el dato $v_{IV}=145{,}15\,\mathrm{m/s}$, el cohete frena de forma uniforme
desde $H=200\,\mathrm{m}$ hasta quedar en reposo ($v=0$) al tocar el suelo:
\[
0=v_{IV}^2-2\,a_{\mathrm{final}}\,H
\;\Rightarrow\;
a_{\mathrm{final}}=\frac{v_{IV}^2}{2H}=\frac{145{,}15^2}{2(200)}\approx 52{,}67\ \mathrm{m/s^2}.
\]
\begin{justif}{Sentido de la aceleraci\'on}
El m\'ovil baja ($-\hat{\jmath}$) y frena, por lo que la aceleraci\'on apunta hacia
arriba, en $+\hat{\jmath}$.
\end{justif}
\begin{resultado}\centering
$\displaystyle \vec{a}_{\mathrm{final}}\approx +52{,}67\,\hat{\jmath}\ \ \mathrm{m/s^2}$
\end{resultado}
\end{paso}

\newpage
% ============================================================
% ============================  P3  ==========================
% ============================================================
\begin{enunciadobox}{Problema 3 --- Desplazamientos del dron \hfill (20 puntos)}
$\vec{A}$: $100\,\mathrm{m}$ a $37^\circ$ al Norte del Este.\quad
$\vec{B}$: $80\,\mathrm{m}$ a $53^\circ$ al Norte del Oeste.\quad
$\vec{C}$: tramo desconocido.\quad
Tras $\vec{A}+\vec{B}+\vec{C}$ el dron queda en
$\vec{D}$: $50\,\mathrm{m}$ al Sur de la base. Ejes: $+x$ Este, $+y$ Norte.
\end{enunciadobox}

% ------------------- P3 (a) -------------------
\begin{paso}{3.a: Componentes de $\vec{A}$, $\vec{B}$ y $\vec{D}$}
\textbf{$\vec{A}$} ($37^\circ$ sobre el eje $+x$, hacia el Norte):
\[
A_x=100\cos 37^\circ\approx 79{,}86,\qquad A_y=100\sin 37^\circ\approx 60{,}18 .
\]
\textbf{$\vec{B}$} ($53^\circ$ desde el Oeste hacia el Norte $\Rightarrow$ $x<0$, $y>0$):
\[
B_x=-80\cos 53^\circ\approx -48{,}15,\qquad B_y=80\sin 53^\circ\approx 63{,}89 .
\]
\textbf{$\vec{D}$} (Sur puro):
\[
D_x=0,\qquad D_y=-50 .
\]
\begin{resultado}\centering
$\vec{A}\approx(79{,}86,\,60{,}18)\,\mathrm{m}\quad
 \vec{B}\approx(-48{,}15,\,63{,}89)\,\mathrm{m}\quad
 \vec{D}=(0,\,-50)\,\mathrm{m}$
\end{resultado}
\end{paso}

% ------------------- P3 (b) -------------------
\begin{paso}{3.b: C\'alculo anal\'itico de $\vec{C}$}
El recorrido total cumple $\vec{A}+\vec{B}+\vec{C}=\vec{D}$, por lo tanto:
\[
\vec{C}=\vec{D}-\vec{A}-\vec{B}.
\]
Componente a componente:
\[
C_x=0-79{,}86-(-48{,}15)\approx -31{,}72,
\]
\[
C_y=-50-60{,}18-63{,}89\approx -174{,}07 .
\]
\begin{resultado}\centering
$\displaystyle \vec{C}\approx(-31{,}72,\,-174{,}07)\ \mathrm{m}$
\end{resultado}
\end{paso}

% ------------------- P3 (c) -------------------
\begin{paso}{3.c: Magnitud y direcci\'on de $\vec{C}$ respecto a $+x$}
\textbf{Magnitud:}
\[
\lVert\vec{C}\rVert=\sqrt{(-31{,}72)^2+(-174{,}07)^2}\approx 176{,}94\ \mathrm{m}.
\]
\textbf{Direcci\'on:} ambas componentes son negativas $\Rightarrow$ tercer cuadrante
(Sur-Oeste). El \'angulo respecto a $+x$ es
\[
\theta=\arctan\!\left(\frac{-174{,}07}{-31{,}72}\right)
\;\Rightarrow\;
\theta\approx -100{,}33^\circ\ \equiv\ 259{,}67^\circ\ (\text{medido CCW desde }+x).
\]
\begin{justif}{Interpretaci\'on}
Equivale a $\approx 79{,}67^\circ$ al Sur del Oeste: el dron debe volver hacia el
suroeste para cerrar el per\'imetro.
\end{justif}
\begin{resultado}\centering
$\displaystyle \lVert\vec{C}\rVert\approx 176{,}94\ \mathrm{m}\qquad
\theta\approx 259{,}67^\circ\ \text{(o } -100{,}33^\circ\text{) respecto a }+x$
\end{resultado}
\end{paso}

% ============================================================
% RESUMEN
% ============================================================
\begin{tcolorbox}[colback=verdePaso!8, colframe=verdePaso, arc=8pt,
    boxrule=1pt, enhanced, breakable,
    title={\bfseries\color{white} Resumen de resultados},
    colbacktitle=verdePaso, fonttitle=\bfseries]
\small
\begin{itemize}[leftmargin=1.6cm, itemsep=3pt]
\item[\textbf{P1 (a)}] $a_1=3{,}00\,\mathrm{m/s^2}$;\quad $a_3=-4{,}00\,\mathrm{m/s^2}$
\item[\textbf{P1 (b)}] $x_1=1{,}5t^2$;\quad $x_2=96+24(t-8)$;\quad $x_3=384+24(t-20)-2(t-20)^2$
\item[\textbf{P1 (c)}] $v_0\approx 14{,}62\,\mathrm{m/s}$;\quad $x^{*}\approx 190{,}09\,\mathrm{m}$ (en $t^{*}\approx 11{,}92\,\mathrm{s}$)
\item[\textbf{P2 (a)}] $y_{\max}\approx 1274{,}89\,\mathrm{m}$
\item[\textbf{P2 (b)}] $\vec{v}_{IV}\approx -145{,}15\,\hat{\jmath}\,\mathrm{m/s}$
\item[\textbf{P2 (c)}] $\vec{a}_{\mathrm{final}}\approx +52{,}67\,\hat{\jmath}\,\mathrm{m/s^2}$
\item[\textbf{P3 (a)}] $\vec{A}\approx(79{,}86,\,60{,}18)$;\ $\vec{B}\approx(-48{,}15,\,63{,}89)$;\ $\vec{D}=(0,\,-50)\,\mathrm{m}$
\item[\textbf{P3 (b)}] $\vec{C}\approx(-31{,}72,\,-174{,}07)\,\mathrm{m}$
\item[\textbf{P3 (c)}] $\lVert\vec{C}\rVert\approx 176{,}94\,\mathrm{m}$;\quad $\theta\approx 259{,}67^\circ\ (-100{,}33^\circ)$
\end{itemize}
\end{tcolorbox}

\vspace{0.8cm}
\begin{center}
    \small\textbf{Usach Premium} --- ¡Nos vemos en la ayudant\'ia! \\
    \textit{"Por y para estudiantes."}
\end{center}

\end{document}
